% SPDX-FileCopyrightText: 2023 SAP SE
%
% SPDX-License-Identifier: Apache-2.0
%
% This file is part of FEDEM - https://openfedem.org

%%%%%%%%%%%%%%%%%%%%%%%%%%%%%%%%%%%%%%%%%%%%%%%%%%%%%%%%%%%%%%%%%%%%%%%%%%%%%%%%
%
% FEDEM Theory Guide.
%
%%%%%%%%%%%%%%%%%%%%%%%%%%%%%%%%%%%%%%%%%%%%%%%%%%%%%%%%%%%%%%%%%%%%%%%%%%%%%%%%

\subsection{The generalized-$\alpha$ method}
\label{subs:The generalized alpha method}

The generalized-$\alpha$ method by Chung and Hulbert~\cite{HulbertChung}
seeks to introduce high frequency dissipation into the numerical solution
by interpolating the inertia forces between time increments $k$ and $k+1$
using a factor $\alpha_m$, and interpolating the elastic, damping and
external forces using another factor $\alpha_f$.
The dynamic equilibrium \eqnref{eq:DynEquil} then takes the form
%
\begin{equation}
\eqalign{
(1-\alpha_m)\Fi_\kp1 + \alpha_m\Fi_k & \;\;+ \cr
(1-\alpha_f)\Fd_\kp1 + \alpha_f\Fd_k & \;\;+ \cr
(1-\alpha_f)\Fs_\kp1 + \alpha_f\Fs_k & \;=\;
(1-\alpha_f)\Q_\kp1  + \alpha_f\Q_k
}\label{eq:GenAlphaEquil}
\end{equation}

The Newmark integration of \eqnref{eq:GenAlphaEquil} is a second order accurate
algorithm provided that
%
\begin{equation}
\gamma = \frac{1}{2} - \alpha_m + \alpha_f
\end{equation}
%
This algorithm is unconditionally stable if the following conditions are met
%
\begin{equation}
\alpha_m \leq \alpha_f \leq \frac{1}{2} \quad\text{and}\quad
\beta \geq \frac{1}{4} + \frac{1}{2}(\alpha_f - \alpha_m)
\end{equation}
%
This region of stability is depicted by the shaded area in
Figure~\ref{fig:GAstab}, which is bounded by the two lines
$\lambda_{1,2}^\infty=-1$ corresponding to $\alpha_m\le\alpha_f$, and
$\lambda_3^\infty=-1$ corresponding to $\alpha_f\le\frac{1}{2}$.
According to~\cite{HulbertChung}, the high frequency modes (relative to the time
increment length $h$) will be damped out if we choose
%
\begin{equation}
\beta = \frac{1}{2}(1 - \alpha_m + \alpha_f)^2
\end{equation}

\begin{figure}[b]
\begin{center}
\setlength{\unitlength}{0.8mm}
\begin{picture}(130,125)(-70,-60)
% Make shaded triangles (the units are [cm])
\put(-50,0){\begin{tikzpicture}
\fill[black!20!white](0,0)--(4,0)--(6,2)--(0,0);
\end{tikzpicture}}
\put(-60,-60){\begin{tikzpicture}
\filldraw[left color=white,right color=black!20!white,draw=white]
(0,0)--(4.8,4.8)--(0.8,4.8)--(0,4.73)--(0,0);
\end{tikzpicture}}
% Axis system
\thinlines
\put(-60,0){\vector(1,0){120}}
\put(0,-60){\vector(0,1){120}}
\multiput(-50,-2)(25,0){5}{\line(0,1){4}}
\multiput(-2,-50)(0,25){5}{\line(1,0){4}}
\put(62,-1){$\alpha_m$}
\put(-1,61){$\alpha_f$}
\put(-55,-8){$-1$}
\put(-30,-8){$-\frac{1}{2}$}
\put(23.7,-8){$\frac{1}{2}$}
\put(48.7,-8){$1$}
\put(3,-51.5){$-1$}
\put(3,-26.5){$-\frac{1}{2}$}
\put(1, 27){$\frac{1}{2}$}
\put(3, 49){$1$}
\put(25,25){\line(-1,0){90}}
%
\thicklines
\put(25,25){\line(-1,-1){85}}
\put(25,25){\line(-3,-1){97}}
\put(0,0){\line(0,1){16.7}}
\put(-50,0){\line(1,0){50}}
\put(0,0){\circle*{1.5}}
\put(25,25){\circle*{1.5}}
\put(0,16.7){\circle*{1.5}}
\put(-50,0){\circle*{1.5}}
\put(-46,-50){$\lambda_{1,2}^\infty=-1$}
\put(-55,26.5){$\lambda_3^\infty=-1$}
\put(-82,-12){$\lambda_3^\infty=\lambda_{1,2}^\infty$}
%
\thinlines
\put(25,9){\vector(-1,0){24.5}}
\put(26,7){HHT-$\alpha$}
\put(-21,14){\vector(0,-1){4}}
\qbezier(25,40)(-21,30)(-21,14)
\put(26,39){Generalized-$\alpha$}
\end{picture}
\end{center}
\caption{Classification of the generalized-$\alpha$ method in the
$\alpha_m-\alpha_f$ space. From~\cite{HulbertChung}.}
\label{fig:GAstab}
\end{figure}

Inserting the internal force linearizations from \meqsref{eq:37}{eq:39} into
\eqnref{eq:GenAlphaEquil} and using $\Q_\kp1=\Q_k+\Delta\Q_k$,
we obtain the generalized-$\alpha$ equation on incremental form
%
\begin{equation}
\eqalign{
(1-\alpha_m){\mf M}_k\Delta\ddot{\mf r}_k + \Fi_k & \;+ \cr
(1-\alpha_f){\mf C}_k\Delta \dot{\mf r}_k + \Fd_k & \;+ \cr
(1-\alpha_f){\mf K}_k\Delta     {\mf r}_k + \Fs_k & \;=\;
(1-\alpha_f)\Delta\Q_k + \Q_k }
\end{equation}
%
or,
by collecting all known quantities of time increment $k$ on the right hand side
%
\begin{equation}
\eqalign{
(1-\alpha_m){\mf M}_k\Delta\ddot{\mf r}_k \;+\; &
(1-\alpha_f){\mf C}_k\Delta \dot{\mf r}_k \;+\;
(1-\alpha_f){\mf K}_k\Delta     {\mf r}_k \cr =\; &
(1-\alpha_f)\Delta\Q_k - \left(\Fi_k + \Fd_k + \Fs_k - \Q_k \right)
}\label{eq:IncGenAlphaEquil}
\end{equation}
%
The last term on the right hand side of \eqnref{eq:IncGenAlphaEquil} can be
recognized as the force residual of time increment $k$, and as such can be
omitted since this should be equal to zero if convergence has been achieved
before advancing to time increment $k+1$.

Inserting the Newmark incremental acceleration and velocity expressions from
\eqsref{eq:322}{eq:324} into \eqnref{eq:IncGenAlphaEquil} yields
%
\begin{equation}
\eqalign{
(1-\alpha_m){\mf M}_k&\left(
\frac{1}{\beta h^2} \Delta{\mf r}_k - {\mf a}_k - \ddot{\mf r}_k \right) \;+ \cr
(1-\alpha_f){\mf C}_k&\left(
\frac{\gamma}{\beta h} \Delta{\mf r}_k - {\mf d}_k - \dot{\mf r}_k \right) \;+ \cr
(1-\alpha_f){\mf K}_k&\Delta{\mf r}_k \;=\;
(1-\alpha_f)\Delta\Q_k - \left(\Fi_k + \Fd_k + \Fs_k - \Q_k \right) }
\end{equation}
%
which after rearranging the known quantities on the right hand side becomes
%
\begin{equation}
{\mf N}_k \Delta{\mf r}_k \;=\; \widehat{\mf R}_k
\end{equation}
%
with
%
\begin{equation}
{\mf N}_k \;=\;
\frac{1-\alpha_m}{\beta h^2}{\mf M}_k \;+\;
\frac{(1-\alpha_f)\gamma}{\beta h} {\mf C}_k \;+\;
(1-\alpha_f){\mf K}_k
\label{eq:NewtonGenAlpha}
\end{equation}
%
and
%
\begin{equation}
\eqalign{
\widehat{\mf R}_k \;=\; &
(1-\alpha_f)\Delta\Q_k - \left(\Fi_k + \Fd_k + \Fs_k - \Q_k \right) \cr +\; &
(1-\alpha_m){\mf M}_k({\mf a}_k + \ddot{\mf r}_k) \;+\;
(1-\alpha_f){\mf C}_k({\mf d}_k + \dot{\mf r}_k)
}\label{eq:PredictorGenAlpha}
\end{equation}

For a linear system, the state at time $t_\kp1$ given by $\Delta{\mf r}_k$,
$\Delta\dot{\mf r}_k$ and $\Delta\ddot{\mf r}_k$, and the associated inertia,
damping and elastic forces forces $\Fi_\kp1$, $\Fd_\kp1$ and $\Fs_\kp1$,
respectively, will satisfy the equilibrium \eqnref{eq:GenAlphaEquil}.
For nonlinear systems, the increments $\Delta{\mf r}_k$, $\Delta\dot{\mf r}_k$
and $\Delta\ddot{\mf r}_k$ will in general not satisfy
\eqnref{eq:GenAlphaEquil}.
In order to ensure dynamic equilibrium before advancing to the next time
increment, the dynamic residual seeks to be minimized through a similar
Newton--Raphson procedure as described in
Section~\ref{s:Newton-Raphson iteration}.
Thus, in each iteration we have to solve the linearized system of equations
%
\begin{equation}
\ls{i}{\mf N}_k \ls{i}\dr_k \;=\; \ls{i-1}{\mf R}^\alpha_k
\end{equation}
%
where the Newton matrix $\ls{i}{\mf N}_k$ is the same as in
\eqnref{eq:NewtonGenAlpha}, and the right-hand-side vector
$\ls{i-1}{\mf R}^\alpha_k$ equals
%
\begin{equation}
\eqalign{
\ls{i-1}{\mf R}^\alpha_k \;=\; &
(1-\alpha_f)\:\ls{i-1}\Q_\kp1  + \alpha_f\Q_k \cr -\; &
(1-\alpha_m)\,\ls{i-1}\Fi_\kp1 - \alpha_m\Fi_k \cr -\; &
(1-\alpha_f)\:\ls{i-1}\Fd_\kp1 - \alpha_f\Fd_k \cr -\; &
(1-\alpha_f)\:\ls{i-1}\Fs_\kp1 - \alpha_f\Fs_k
}\label{eq:CorrectorGenAlpha}
\end{equation}
